\section{实验与结果分析}

\subsection{实验环境与数据集}

【列出硬件环境、软件环境与所用数据集的基本信息。】

实验硬件环境为……，软件环境为……。实验采用……数据集，该数据集包含……条样本，按 …… 的比例划分为训练集、验证集与测试集。

\subsection{实验结果与分析}

【给出对比实验与消融实验的结果表，并对表中数据进行分析，说明各模块的贡献。】

为验证所提方法的有效性，本文选取……作为对比方法，评价指标采用……。实验结果如表 \ref{tab:result} 所示。

\begin{table}[htbp]
  \centering
  \caption{对比实验结果（示例，请替换为本人论文的实验结果）}
  \label{tab:result}
  \vspace{0.8em}
  \renewcommand{\arraystretch}{1.3}
  \begin{tabular}{lccc}
    \toprule
    \textbf{方法} & \textbf{指标一} & \textbf{指标二} & \textbf{指标三} \\
    \midrule
    基线方法 A & 0.0000 & 0.0000 & 0.0000 \\
    基线方法 B & 0.0000 & 0.0000 & 0.0000 \\
    本文方法   & \textbf{0.0000} & \textbf{0.0000} & \textbf{0.0000} \\
    \bottomrule
  \end{tabular}
\end{table}

由表 \ref{tab:result} 可见，……。为进一步验证各模块的作用，本文还进行了消融实验，结果显示……。

\subsection{本章小结}

本章介绍了实验设置，并通过对比实验与消融实验验证了所提方法的有效性。
